% Original knowledge-map reconstruction from Propositional Logic slides 63--68.
% Unlike the earlier lecture's (x,y), subscripts here are row,column.
\begin{tikzpicture}[x=1.55cm,y=1.15cm,every node/.style={font=\small}]
  \fill[noteBlue!9] (0,0) rectangle (2,1);
  \fill[noteBlue!9] (0,1) rectangle (1,2);
  \fill[verifyOrange!16] (0,2) rectangle (1,3);
  \draw[step=1,black!50] (0,0) grid (3,3);
  \foreach \k in {1,2,3} {
    \node at (\k-0.5,-0.25) {\k};
    \node at (-0.25,\k-0.5) {\k};
  }
  \node at (1.5,-0.63) {列 $c$（横向）};
  \node[rotate=90] at (-0.67,1.5) {行 $r$（纵向）};
  \node[align=center] at (0.5,0.5) {$\neg S_{11}$\\$\neg W_{11}$};
  \node[align=center] at (1.5,0.5) {$\neg S_{12}$\\$\neg W_{12}$};
  \node at (2.5,0.5) {$\neg W_{13}$};
  \node[align=center] at (0.5,1.5) {$A_{21},S_{21}$\\$\neg W_{21}$};
  \node at (1.5,1.5) {$\neg W_{22}$};
  \node[text=verifyOrange,align=center] at (0.5,2.5) {$W_{31}$\\已确定};
  \draw[-{Latex},noteBlue,thick] (3.5,0.5) -- (3.5,1.3);
  \node[anchor=south] at (3.5,1.3) {North};
\end{tikzpicture}
